\documentclass[a4paper,11pt]{article}
\usepackage[utf8]{inputenc}
\usepackage[spanish]{babel}
\usepackage{geometry}
\usepackage{enumitem}
\usepackage{sectsty}
\usepackage{titlesec}
\usepackage{tikz}
\usepackage{array}

\geometry{
    a4paper,
    left=20mm,
    top=18mm,
    right=20mm,
    bottom=20mm
}

% Sin números de página
\pagestyle{empty}

% Estilo de secciones con números en círculos
\newcommand{\circled}[1]{%
    \tikz[baseline=(char.base)]{%
        \node[shape=circle,draw,inner sep=0.5pt,fill=black,text=white,font=\bfseries\footnotesize] (char) {#1};%
    }%
}

% Formato de secciones
\titleformat{\section}
{\normalfont\Large\bfseries}
{\circled{\thesection}}{0.5em}{}

\titlespacing*{\section}{0pt}{12pt}{8pt}

\titleformat{\subsection}
{\normalfont\normalsize\bfseries}
{}{0em}{}

\titlespacing*{\subsection}{8pt}{4pt}{4pt}

\begin{document}

% Encabezado del curso
\begin{center}
\small{Universidad de Granada}

\vspace{0.3cm}

\textbf{\large{Producción e Interacción Oral - Nivel 7}}

\vspace{0.4cm}

\small
\begin{tabular}{@{}ll@{}}
\textbf{Grupos:} & A (Aula 24) y B (Aula 08) \\
\textbf{Centro:} & Centro de Lenguas Modernas / Edificio A \\
\textbf{Profesor:} & Javier Benítez Lainez \\
\textbf{Tutorías:} & Martes 9:00-10:00 y 14:30-15:30 \\
& Sala de profesores / Email: benitezl@go.ugr.es
\end{tabular}
\end{center}

\vspace{0.5cm}
\section*{Ejercicios Prácticos de Contraargumentación}

\vspace{0.4cm}

\section{Cuaderno de Actividades para Debate C1}

\vspace{0.2cm}

---

\vspace{0.2cm}

\section{INTRODUCCIÓN}

\vspace{0.2cm}

\subsection{Propósito de este cuaderno}

Este documento contiene ejercicios prácticos para desarrollar y afianzar tus habilidades de contraargumentación en español nivel C1. Los ejercicios están organizados por dificultad y tipo de habilidad.

\vspace{0.2cm}

\subsection{Cómo usar este cuaderno}

1. \textbf{Lee el argumento} presentado

2. \textbf{Prepara tu refutación} usando las estructuras aprendidas

3. \textbf{Practica en voz alta} cronometrando tu respuesta

4. \textbf{Evalúate} usando las rúbricas proporcionadas

5. \textbf{Repite} intentando mejorar cada vez

\vspace{0.2cm}

---

\vspace{0.2cm}

\section{PARTE 1: EJERCICIOS DE RECONOCIMIENTO Y CONCESIÓN}

\vspace{0.2cm}

\subsection{Ejercicio 1.1: Practicar el reconocimiento}

\vspace{0.2cm}

\textbf{Instrucciones:} Para cada argumento oponente, escribe una frase de reconocimiento apropiada.

\vspace{0.2cm}

\textbf{Argumento 1:}

"La tecnología está destruyendo las relaciones interpersonales porque la gente pasa más tiempo en sus teléfonos que hablando cara a cara."

\vspace{0.2cm}

Tu reconocimiento: _________________________

\vspace{0.2cm}

\textbf{Argumento 2:}

"El gobierno no debería financiar las artes porque es un gasto innecesario de dinero público."

\vspace{0.2cm}

Tu reconocimiento: _________________________

\vspace{0.2cm}

\textbf{Argumento 3:}

"El matrimonio es una institución obsoleta que ya no tiene relevancia en la sociedad moderna."

\vspace{0.2cm}

Tu reconocimiento: _________________________

\vspace{0.2cm}

\textbf{Argumento 4:}

"Los videojuegos violentos causan comportamiento violento en los jóvenes."

\vspace{0.2cm}

Tu reconocimiento: _________________________

\vspace{0.2cm}

\textbf{Argumento 5:}

"El comercio internacional es perjudicial porque destruye empleos locales."

\vspace{0.2cm}

Tu reconocimiento: _________________________

\vspace{0.2cm}

---

\vspace{0.2cm}

\subsection{Ejercicio 1.2: Practicar la concesión parcial}

\vspace{0.2cm}

\textbf{Instrucciones:} Para cada argumento, escribe una concesión que reconozca validez parcial.

\vspace{0.2cm}

\textbf{Argumento 1:}

"La inteligencia artificial va a reemplazar a todos los trabajadores humanos."

\vspace{0.2cm}

Concesión: _________________________

\vspace{0.2cm}

\textbf{Argumento 2:}

"Todos los políticos son corruptos y no sirven para nada."

\vspace{0.2cm}

Concesión: _________________________

\vspace{0.2cm}

\textbf{Argumento 3:}

"El cambio climático es un hoax inventado por científicos para obtener financiación."

\vspace{0.2cm}

Concesión: _________________________

\vspace{0.2cm}

---

\vspace{0.2cm}

\section{PARTE 2: EJERCICIOS DE REFUTACIÓN ESPECÍFICA}

\vspace{0.2cm}

\subsection{Ejercicio 2.1: Refutación por evidencia contraria}

\vspace{0.2cm}

\textbf{Instrucciones:} Refuta cada argumento presentando evidencia contraria específica.

\vspace{0.2cm}

\textbf{Argumento:}

"Las mujeres no son aptas para liderazgo en negocios porque son demasiado emocionales."

\vspace{0.2cm}

Tu refutación con evidencia:

\begin{itemize}

  \item Reconocimiento: _________________________

  \item Evidencia contraria: _________________________

  \item Conexión: _________________________

\end{itemize}

\vspace{0.2cm}

\textbf{Argumento:}

"La inmigración causa aumento de criminalidad."

\vspace{0.2cm}

Tu refutación con evidencia:

\begin{itemize}

  \item Reconocimiento: _________________________

  \item Evidencia contraria: _________________________

  \item Conexión: _________________________

\end{itemize}

\vspace{0.2cm}

---

\vspace{0.2cm}

\subsection{Ejercicio 2.2: Refutación por fallo lógico}

\vspace{0.2cm}

\textbf{Instrucciones:} Identifica el fallo lógico y refuta el argumento.

\vspace{0.2cm}

\textbf{Argumento:}

"No deberíamos tener leyes de control de armas porque las personas que cometen crímenes no siguen las leyes de todos modos."

\vspace{0.2cm}

Fallo lógico: _________________________

\vspace{0.2cm}

Tu refutación: _________________________

\vspace{0.2cm}

\textbf{Argumento:}

"Si permitimos el matrimonio gay, pronto la gente querrá casarse con sus perros."

\vspace{0.2cm}

Fallo lógico: _________________________

\vspace{0.2cm}

Tu refutación: _________________________

\vspace{0.2cm}

---

\vspace{0.2cm}

\subsection{Ejercicio 2.3: Refutación por analogía}

\vspace{0.2cm}

\textbf{Instrucciones:} Crea una analogía para refutar cada argumento.

\vspace{0.2cm}

\textbf{Argumento:}

"No deberíamos tener impuestos porque el gobierno simplemente desperdicia el dinero."

\vspace{0.2cm}

Analogía: _________________________

\vspace{0.2cm}

Refutación completa: _________________________

\vspace{0.2cm}

\textbf{Argumento:}

"Los estudiantes no deber tener que asistir a clases; pueden aprender todo por internet."

\vspace{0.2cm}

Analogía: _________________________

\vspace{0.2cm}

Refutación completa: _________________________

\vspace{0.2cm}

---

\vspace{0.2cm}

\subsection{Ejercicio 2.4: Refutación por distinción}

\vspace{0.2cm}

\textbf{Instrucciones:} Identifica una distinción importante y úsala para refutar.

\vspace{0.2cm}

\textbf{Argumento:}

"La eutanasia es lo mismo que asesinato."

\vspace{0.2cm}

Distinción: _________________________

\vspace{0.2cm}

Refutación: _________________________

\vspace{0.2cm}

\textbf{Argumento:}

"Críticas al gobierno son lo mismo que traición."

\vspace{0.2cm}

Distinción: _________________________

\vspace{0.2cm}

Refutación: _________________________

\vspace{0.2cm}

---

\vspace{0.2cm}

\section{PARTE 3: EJERCICIOS DE REFUTACIÓN COMPLETA (RE-C-R)}

\vspace{0.2cm}

\subsection{Ejercicio 3.1: Temas sociales}

\vspace{0.2cm}

\textbf{Instrucciones:} Prepara una refutación completa usando el modelo RE-C-R. Cronometra: 60 segundos.

\vspace{0.2cm}

\textbf{Argumento 1:}

"La justicia restaurativa no funciona. Los criminales solo entienden el castigo duro."

\vspace{0.2cm}

R - Reconocer: _________________________

\vspace{0.2cm}

C - Conceder (si aplica): _________________________

\vspace{0.2cm}

R - Refutar: _________________________

\vspace{0.2cm}

C - Conectar: _________________________

\vspace{0.2cm}

---

\vspace{0.2cm}

\textbf{Argumento 2:}

"El movimiento feminista ya no es necesario porque las mujeres y los hombres son iguales ahora."

\vspace{0.2cm}

R - Reconocer: _________________________

\vspace{0.2cm}

C - Conceder (si aplica): _________________________

\vspace{0.2cm}

R - Refutar: _________________________

\vspace{0.2cm}

C - Conectar: _________________________

\vspace{0.2cm}

---

\vspace{0.2cm}

\subsection{Ejercicio 3.2: Temas económicos}

\vspace{0.2cm}

\textbf{Argumento 1:}

"El salario mínimo causa desempleo porque las empresas no pueden pagar salarios altos."

\vspace{0.2cm}

R - Reconocer: _________________________

\vspace{0.2cm}

C - Conceder (si aplica): _________________________

\vspace{0.2cm}

R - Refutar: _________________________

\vspace{0.2cm}

C - Conectar: _________________________

\vspace{0.2cm}

---

\vspace{0.2cm}

\textbf{Argumento 2:}

"La globalización es terrible porque destruye la cultura local."

\vspace{0.2cm}

R - Reconocer: _________________________

\vspace{0.2cm}

C - Conceder (si aplica): _________________________

\vspace{0.2cm}

R - Refutar: _________________________

\vspace{0.2cm}

C - Conectar: _________________________

\vspace{0.2cm}

---

\vspace{0.2cm}

\subsection{Ejercicio 3.3: Temas ambientales}

\vspace{0.2cm}

\textbf{Argumento 1:}

"Las acciones individuales contra el cambio climático no sirven para nada. Solo importa lo que hacen los gobiernos y las grandes empresas."

\vspace{0.2cm}

R - Reconocer: _________________________

\vspace{0.2cm}

C - Conceder (si aplica): _________________________

\vspace{0.2cm}

R - Refutar: _________________________

\vspace{0.2cm}

C - Conectar: _________________________

\vspace{0.2cm}

---

\vspace{0.2cm}

\textbf{Argumento 2:}

"La energía nuclear es demasiado peligrosa y nunca deberíamos usarla."

\vspace{0.2cm}

R - Reconocer: _________________________

\vspace{0.2cm}

C - Conceder (si aplica): _________________________

\vspace{0.2cm}

R - Refutar: _________________________

\vspace{0.2cm}

C - Conectar: _________________________

\vspace{0.2cm}

---

\vspace{0.2cm}

\subsection{Ejercicio 3.4: Temas tecnológicos}

\vspace{0.2cm}

\textbf{Argumento 1:}

"La privacidad es una cosa del pasado. No podemos hacer nada para protegerla en la era digital, así que no importa."

\vspace{0.2cm}

R - Reconocer: _________________________

\vspace{0.2cm}

C - Conceder (si aplica): _________________________

\vspace{0.2cm}

R - Refutar: _________________________

\vspace{0.2cm}

C - Conectar: _________________________

\vspace{0.2cm}

---

\vspace{0.2cm}

\textbf{Argumento 2:}

"Los algoritmos de redes sociales deberían ser prohibidos porque causan adicción y polarización."

\vspace{0.2cm}

R - Reconocer: _________________________

\vspace{0.2cm}

C - Conceder (si aplica): _________________________

\vspace{0.2cm}

R - Refutar: _________________________

\vspace{0.2cm}

C - Conectar: _________________________

\vspace{0.2cm}

---

\vspace{0.2cm}

\section{PARTE 4: EJERCICIOS DE RESPUESTA RÁPIDA}

\vspace{0.2cm}

\subsection{Ejercicio 4.1: Refutaciones de 30 segundos}

\vspace{0.2cm}

\textbf{Instrucciones:} Prepara una refutación concisa de máximo 30 segundos para cada argumento.

\vspace{0.2cm}

\textbf{1. "El dinero compra la felicidad."}

Tu refutación (30 seg): _________________________

\vspace{0.2cm}

\textbf{2. "Todo el mundo tiene las mismas oportunidades en la vida."}

Tu refutación (30 seg): _________________________

\vspace{0.2cm}

\textbf{3. "La experiencia es más importante que la educación formal."}

Tu refutación (30 seg): _________________________

\vspace{0.2cm}

\textbf{4. "La gente es pobre porque es perezosa."}

Tu refutación (30 seg): _________________________

\vspace{0.2cm}

\textbf{5. "La pena de muerte previene el crimen."}

Tu refutación (30 seg): _________________________

\vspace{0.2cm}

---

\vspace{0.2cm}

\subsection{Ejercicio 4.2: Refutaciones de 15 segundos (brevidad extrema)}

\vspace{0.2cm}

\textbf{Instrucciones:} Refuta en MAXIMO 15 segundos siendo aún efectivo.

\vspace{0.2cm}

\textbf{1. "La gente buena no necesita leyes para comportarse bien."}

Tu refutación (15 seg): _________________________

\vspace{0.2cm}

\textbf{2. "Si no tienes nada que esconder, no deberías preocuparte por la vigilancia."}

Tu refutación (15 seg): _________________________

\vspace{0.2cm}

\textbf{3. "El mercado libre soluciona todos los problemas económicos."}

Tu refutación (15 seg): _________________________

\vspace{0.2cm}

---

\vspace{0.2cm}

\section{PARTE 5: ESCENARIOS DE DEBATE COMPLETO}

\vspace{0.2cm}

\subsection{Ejercicio 5.1: Debate estructurado}

\vspace{0.2cm}

\textbf{Tema:} "¿Deberían los gobiernos prohibir la publicidad de alimentos no saludables?"

\vspace{0.2cm}

\textbf{Posición A (Prohibición):}

Prepara 3 argumentos a favor.

\vspace{0.2cm}

Argumento 1: _________________________

\vspace{0.2cm}

Argumento 2: _________________________

\vspace{0.2cm}

Argumento 3: _________________________

\vspace{0.2cm}

\textbf{Posición B (En contra):}

Prepara refutaciones para cada argumento de la Posición A.

\vspace{0.2cm}

Refutación al Argumento 1: _________________________

\vspace{0.2cm}

Refutación al Argumento 2: _________________________

\vspace{0.2cm}

Refutación al Argumento 3: _________________________

\vspace{0.2cm}

---

\vspace{0.2cm}

\subsection{Ejercicio 5.2: Debate público simulado}

\vspace{0.2cm}

\textbf{Tema:} "¿Debería la educación universitaria ser gratuita para todos?"

\vspace{0.2cm}

\textbf{Rol 1: Primer orador (Pro-gratuidad)}

Tienes 2 minutos para presentar tus argumentos principales.

\vspace{0.2cm}

Tu presentación: _________________________

\vspace{0.2cm}

\textbf{Rol 2: Segundo orador (Anti-gratuidad)}

Tienes 2 minutos para refutar los argumentos del Rol 1 y presentar los tuyos.

\vspace{0.2cm}

Tu presentación: _________________________

\vspace{0.2cm}

\textbf{Rol 3: Réplica (Pro-gratuidad)}

Tienes 1 minuto para contraargumentar las críticas del Rol 2.

\vspace{0.2cm}

Tu réplica: _________________________

\vspace{0.2cm}

---

\vspace{0.2cm}

\section{PARTE 6: ANÁLISIS DE REFUTACIONES}

\vspace{0.2cm}

\subsection{Ejercicio 6.1: Analiza refutaciones modelo}

\vspace{0.2cm}

\textbf{Instrucciones:} Lee estas refutaciones y analiza qué estrategias usan.

\vspace{0.2cm}

\textbf{Ejemplo 1:}

\vspace{0.2cm}

\textbf{Argumento oponente:} "La pena de muerte es necesaria para disuadir el crimen."

\vspace{0.2cm}

\textbf{Refutación:}

"El oponente sostiene que la pena de muerte disuade el crimen. Si bien es comprensible el deseo de prevenir crímenes, la evidencia no apoya esta posición. De hecho, estados de EE.UU. con pena de muerte tienen tasas de criminalidad similares o más altas que estados sin ella. Además, países que abolieron la pena de muerte no han visto aumentos de criminalidad. Por consiguiente, la pena de muerte no es efectiva como disuasión."

\vspace{0.2cm}

\textbf{Análisis:}

\begin{itemize}

  \item ¿Reconoce el argumento? _________

  \item ¿Concede algo? _________

  \item ¿Qué tipo de evidencia presenta? _________

  \item ¿Cómo conecta de vuelta? _________

\end{itemize}

\vspace{0.2cm}

---

\vspace{0.2cm}

\textbf{Ejemplo 2:}

\vspace{0.2cm}

\textbf{Argumento oponente:} "Deberíamos prohibir todas las armas de fuego."

\vspace{0.2cm}

\textbf{Refutación:}

"Entiendo la preocupación sobre la violencia armada. Es cierto que este es un problema serio. Sin embargo, una prohibición total es como prohibir todos los coches porque hay accidentes. La solución no es prohibición absoluta sino regulación inteligente: licencias, verificación de antecedentes, y capacitación obligatoria funcionan mejor que prohibiciones totales, como demuestra el ejemplo de Suiza y Canadá."

\vspace{0.2cm}

\textbf{Análisis:}

\begin{itemize}

  \item ¿Usa analogía? _________

  \item ¿Qué propone como alternativa? _________

  \item ¿Qué ejemplos específicos cita? _________

\end{itemize}

\vspace{0.2cm}

---

\vspace{0.2cm}

\section{PARTE 7: AUTOEVALUACIÓN}

\vspace{0.2cm}

\subsection{Lista de control para refutación}

\vspace{0.2cm}

\textbf{Después de cada ejercicio, evalúate:}

\vspace{0.2cm}

\section{Estructura}

\begin{itemize}

  \item [ ] Reconocí el argumento oponente

  \item [ ] Concedí validez parcial cuando aplicaba

  \item [ ] Presenté refutación clara

  \item [ ] Conecté de vuelta a mi posición

\end{itemize}

\vspace{0.2cm}

\section{Contenido}

\begin{itemize}

  \item [ ] Usé evidencia específica

  \item [ ] Evité falacias lógicas

  \item [ ] Mi lógica fue clara

  \item [ ] Cité fuentes o ejemplos

\end{itemize}

\vspace{0.2cm}

\section{Expresión}

\begin{itemize}

  \item [ ] Usé conectores apropiados

  \item [ ] Mi vocabulario fue preciso

  \item [ ] Mi gramática fue correcta

  \item [ ] Mantuve un tono respetuoso

\end{itemize}

\vspace{0.2cm}

\section{Tiempo}

\begin{itemize}

  \item [ ] Me ajusté al tiempo límite

  \item [ ] No hablé demasiado rápido

  \item [ ] No me quedé corto

\end{itemize}

\vspace{0.2cm}

---

\vspace{0.2cm}

\section{PARTE 8: EJERCICIOS CON PAREJA}

\vspace{0.2cm}

\subsection{Ejercicio 8.1: Refutación en vivo}

\vspace{0.2cm}

\textbf{Instrucciones:} En parejas, uno presenta un argumento y el otro debe refutarlo inmediatamente.

\vspace{0.2cm}

\textbf{Pareja 1:}

\begin{itemize}

  \item Persona A: Argumenta a favor de X

  \item Persona B: Refuta inmediatamente (RE-C-R completo)

\end{itemize}

\vspace{0.2cm}

\textbf{Pareja 2:}

\begin{itemize}

  \item Persona A: Presenta un nuevo argumento

  \item Persona B: Refuta inmediatamente

\end{itemize}

\vspace{0.2cm}

\textbf{Después:} Intercambien roles. Persona B argumenta, Persona A refuta.

\vspace{0.2cm}

---

\vspace{0.2cm}

\subsection{Ejercicio 8.2: Debate de ping-pong}

\vspace{0.2cm}

\textbf{Instrucciones:} Alterna argumentos y contraargumentos sin preparación.

\vspace{0.2cm}

\begin{itemize}

  \item Persona A: Argumento inicial (30 segundos)

  \item Persona B: Refutación (30 segundos)

  \item Persona A: Contra-refutación (30 segundos)

  \item Persona B: Segunda refutación (30 segundos)

\end{itemize}

\vspace{0.2cm}

\textbf{Tema a elegir:} _________________________

\vspace{0.2cm}

---

\vspace{0.2cm}

\section{SOLUCIONES SUGERIDAS}

\vspace{0.2cm}

\subsection{Ejercicio 1.1 (Reconocimiento)}

1. "El oponente argumenta que la tecnología está destruyendo relaciones interpersonales debido al tiempo excesivo en dispositivos."

2. "El equipo contrario sostiene que el gobierno no debería financiar las artes por ser un gasto innecesario."

3. "Mi interlocutor afirma que el matrimonio es una institución obsoleta sin relevancia moderna."

4. "Se sostiene que los videojuegos violentos causan comportamiento violento juvenil."

5. "El argumento es que el comercio internacional es perjudicial porque destruye empleos locales."

\vspace{0.2cm}

\subsection{Ejercicio 2.2 (Fallo lógico)}

1. \textbf{Fallo:} Falacia de hombre de paja + error lógico sobre la función de las leyes.

\textbf{Refutación:} "Ese argumento ignora que las leyes no solo castigan sino también previenen y establecen normas sociales. Además, regula el acceso: facinerosos obtienen armas más fácilmente sin leyes."

\vspace{0.2cm}

2. \textbf{Fallo:} Falacia de pendiente resbaladiza.

\textbf{Refutación:} "Esto es una generalización excesiva. El matrimonio entre adultos consentidores es completamente diferente de relaciones no consentientes. No hay conexión lógica entre ambos."

\vspace{0.2cm}

---

\vspace{0.2cm}

\textbf{Sesión 6: Contraargumentación - Ejercicios Prácticos}

\textbf{Nivel C1 - Español Oral Avanzado}

\vspace{0.2cm}

\end{document}
